% Logical Methods in Computer Science.  With lmcs.cls on the TeX path this builds
% in the LMCS class; otherwise it falls back to article so the draft still builds.
\IfFileExists{lmcs.cls}{%
  \documentclass{lmcs}\def\lmcsclass{}}{%
  \documentclass[11pt]{article}\usepackage[margin=1.1in]{geometry}}
\usepackage{amsmath,amssymb,amsthm}
\usepackage[numbers]{natbib}
\usepackage{booktabs}
\usepackage{microtype}
\ifdefined\lmcsclass\else
  \usepackage{hyperref}
  \providecommand{\keywords}[1]{\par\noindent\textbf{Keywords.} #1\par}
\fi
\ifdefined\lmcsclass
  % LMCS predefines its theorem environments; ours are aliases, sharing its numbering.
  \newenvironment{theorem}{\thm}{\endthm}
  \newenvironment{lemma}{\lem}{\endlem}
  \newenvironment{proposition}{\prop}{\endprop}
  \newenvironment{corollary}{\cor}{\endcor}
  \newenvironment{definition}{\defi}{\enddefi}
  \newenvironment{remark}{\rem}{\endrem}
\else
  \theoremstyle{plain}
  \newtheorem{theorem}{Theorem}[section]
  \newtheorem{lemma}[theorem]{Lemma}
  \newtheorem{proposition}[theorem]{Proposition}
  \newtheorem{corollary}[theorem]{Corollary}
  \theoremstyle{definition}
  \newtheorem{definition}[theorem]{Definition}
  \theoremstyle{remark}
  \newtheorem{remark}[theorem]{Remark}
\fi
\newcommand{\code}[1]{\langle #1\rangle}
\newcommand{\red}{\twoheadrightarrow}
\newcommand{\hred}{\twoheadrightarrow_{h}}
\newcommand{\lean}[1]{\texttt{#1}}
\newcommand{\fn}[1]{\textup{#1}}  % named terms: upright everywhere, even in theorem statements
\newcommand{\T}{\mathsf{T}}
\newcommand{\F}{\mathsf{F}}
\newcommand{\num}[1]{\underline{#1}}
\newcommand{\Nat}{\mathtt{Nat}}
\newcommand{\Bool}{\mathtt{Bool}}

\newcommand{\theabstract}{%
In weak combinatory logic over $\{S,K,I\}$ no term quotes every term, so a staged
language built on it fixes quotation at compile time.  We read that discipline as the
S4 necessity modality of Davies--Pfenning and Kavvos, and ask where a
\emph{behavioral lift} exists: a closed term that sends a live value to a code
agreeing with it, an instance of Kavvos' Q-coalgebras and a weak substitute for the
missing unit.  The answer depends on where the value must be total.  A lift exists at
$\Bool\to\Bool$ under totality on the booleans; provably none exists at $\Nat\to\Bool$
under totality on the numerals.  Under a modulus, an operational
form of totality on a compact point set, a lift exists at $\Nat\to\Bool$ by search
over $\mathbb{N}_\infty$ and at $(\Nat\to\Bool)\to\Bool$ by Berger's search over
Cantor space.  Totality at the point at infinity suffices for $\Nat\to\Bool$, and
totality on every point of Cantor space suffices in an oracle semantics; totality on
the definable streams does not, by a Kleene tree built inside the calculus.  The staging
line for behavioral lift is therefore drawn at compactness of the point set on which
totality is demanded, namely Escard\'o's boundary for decidable equality.  We construct
the modal reading and produce constructions in weak combinatory logic and a
compactness observation.  All results are mechanized in Lean~4.}

\begin{document}

\ifdefined\lmcsclass
  \title[The staging line is compactness]{The Staging Line Is
    Compactness: Behavioral Reification in Weak Combinatory Logic}
  \titlecomment{All results are mechanized in Lean~4 (\texttt{lean-ski}).}
  \author[N.~E.~Davis]{N.~E.~Davis}[a]
  \address{ZeroAttest, Champaign, Illinois}
  \keywords{combinatory logic, staging, modal type theory, reification, exhaustible sets,
  fan theorem, Kleene tree, Lean}
  \begin{abstract}\theabstract\end{abstract}
  \maketitle
\else
  \title{The Staging Line Is Compactness\\[2pt]\large Behavioral reification in weak
    combinatory logic\\[2pt]\normalsize(working draft, LMCS format pending \texttt{lmcs.cls})}
  \author{N.\ E.\ Davis}
  \date{}
  \maketitle
  \begin{abstract}\theabstract\end{abstract}
  \keywords{combinatory logic, staging, modal type theory, reification, exhaustible sets,
  fan theorem, Kleene tree, Lean}
\fi


\section{Introduction}\label{sec:intro}
A language whose instruction set is $\{S,K,I\}$ cannot quote at run time: no closed
term $Q$ has $Q\,M \red \code{M}$ for every $M$
\citep{Barendregt1991,DavisMeta} (\lean{no\_quote}).  The declarative $SKI$ ``assembly
language'' SKIjack \citep{skijack}
therefore makes quotation a compile-time construct, and checks that four surface
positions (equality operands, quotation arguments, namespace keys and scry paths)
hold data, not functions.  Davies and Pfenning \citep{DaviesPfenning2001} read
staging as S4 necessity, and Kavvos \citep{Kavvos2017} develops intensional
recursion in that frame for untyped $\lambda$, where evaluation and code-for-code
are definable and only quoting is not.  Section~\ref{sec:modality} places SKIjack's
discipline in the same frame, with the comonad's structure maps mechanized over
$\{S,K,I\}$.

While the syntactic unit is impossible, a weaker one need not be.  A \emph{behavioral}
lift sends a live value to the code of \emph{some} term that behaves as it does.
We call the boundary between the live values lifted to codes by some closed term and those
lifted by no closed term the \emph{staging line}.
Davies and Pfenning remark that a lift $A \to \Box A$ exists exactly when every
arrow of $A$ is under a box \citep[\S5.3]{DaviesPfenning2001}, a criterion suited to
structural copies.  MetaML forbids lifting functions by design
\citep{TahaSheard2000}, and type-directed partial evaluation reaches sums only
with control and stops at inductive types \citep{Danvy1996}.  We show, inside weak
combinatory logic and with mechanized proofs, that the behavioral question has a
different answer, one that turns on where the value must be total
(Table~\ref{tab:results}).
\begin{table}[tb]\centering\footnotesize
\caption{behavioral lifts, by type and by where totality is demanded.}\label{tab:results}
\begin{tabular}{@{}lll@{}}
\toprule
 & totality on a non-compact point set & totality on a compact point set \\
\midrule
$\Bool\to\Bool$ & --- & lift (Thm.~\ref{thm:bool}) \\
$\Nat\to\Bool$ & $\mathbb{N}$: no lift (Thm.~\ref{thm:nolift}) & $\mathbb{N}_\infty$: lift (Thms.~\ref{thm:nat}, \ref{thm:natfan}) \\
$(\Nat\to\Bool)\to\Bool$ & definable streams: no lift (Thm.~\ref{thm:ctotal}) & all points: lift$^\ddagger$ (Thms.~\ref{thm:cantor}, \ref{thm:ofan}) \\
\bottomrule
\end{tabular}\\[2pt]
{\footnotesize $^\ddagger$in an oracle semantics, by the classical fan theorem.}
\end{table}
A function has a \emph{modulus} when it reads its argument only to a bounded
depth: boundedly many successor layers of a numeral, or boundedly many bits of a
stream.  A modulus is an operational form of totality on a compact point set:
it asks for a value on every term that agrees with a point up to a finite depth.
For $\Nat\to\Bool$ the points correspond to $\mathbb{N}_\infty$, the naturals with a
point at infinity; for $(\Nat\to\Bool)\to\Bool$ they are the points of Cantor space,
the infinite boolean sequences $\alpha:\mathbb{N}\to\Bool$, of which terms define only
countably many.
Totality on the numerals does not demand a value at $\infty$, and under that
hypothesis no lift exists (Theorem~\ref{thm:nolift}).

The failure is sharp: one function that is total on every numeral but has no value at
$\infty$, added to the functions with a modulus, defeats every lift
(Theorem~\ref{thm:sharp}).  The obstruction is intensional; extensionally, the
functions with a modulus are eventually constant and the witness is constantly false.
The \emph{staging line} can be read as the compactness of the point set on which
totality is demanded. This recapitulates Escard\'o's boundary between compact and
discrete types for decidable equality \citep[Thm.~4.13]{Escardo2008}, but for code
extraction.  % We offer the reading as an observation, not a classification; what we
% prove is each cell of the table above.
% We prove each cell of Table~\ref{tab:results}.

\paragraph{Contributions.}
\begin{enumerate}
\item The S4 reading of compile-time quotation in $\{S,K,I\}$, with re-encoding and
  intensional fixed points mechanized, and a definition of behavioral lift
  parameterized by a hypothesis and a set of test inputs (\S\ref{sec:staging}).
\item Lifts from a modulus, inside weak combinatory logic: at $\Bool\to\Bool$, at
  $\Nat\to\Bool$, and at Cantor space, each returning a canonical code and each
  idempotent, so that Kavvos' second coalgebra law holds up to conversion and fails
  at his $\sim$ (\S\ref{sec:modulus}).
\item When totality yields a modulus (\S\ref{sec:fan}): on the numerals it does not,
  and no lift exists, by a stuttering simulation; on the points of $\mathbb{N}_\infty$
  it does outright; on the definable streams it does not, by a Kleene tree built
  inside the calculus; and on all points of Cantor space it does, in an oracle
  semantics.
\item Consequences for a staged language's data rule, including a fix to SKIjack's
  checker (\S\ref{sec:skijack}).
\end{enumerate}
Every result is mechanized in Lean~4 against \texttt{lean-ski} \citep{leanski}
(\S\ref{sec:mech}).  The
ingredients are known: finite tabulation \citep{AltenkirchUustalu2004}, Berger's
search \citep{Berger1990} and Escard\'o's account of it
\citep{Escardo2007,Escardo2008}, and the continuity argument behind finite
identification \citep{Balbach2020}.  Section~\ref{sec:related} says what is theirs.

%% ---------------------------------------------------------------------------
\section{Preliminaries}\label{sec:prelim}

\paragraph{Terms and reduction.}  Terms are $S$, $K$, $I$ and applications
$t\,u$, associating to the left; all terms are closed.  Weak reduction $\red$ is the
reflexive--transitive closure of $S\,f\,g\,x \to f\,x\,(g\,x)$, $K\,x\,y \to x$,
$I\,x \to x$ anywhere in a term.  Head reduction $\hred$ contracts only at the head
of the spine.  A term is in weak head normal form when its head combinator has too
few arguments to contract.  We use three standard facts, all mechanized:
Church--Rosser for $\red$ \citep[Thm.~2.15]{HindleySeldin2008} (\lean{church\_rosser});
a form of standardization \citep[\S11B3]{CHS1972}, \citep[Rem.~3.24]{HindleySeldin2008},
by which every weak reduction to a weak head normal form yields a head reduction to
one (\lean{eval\_back}, \lean{Compose.eval\_K\_iff}); and determinism of head
reduction, which is immediate because a term has at most one head redex.

\paragraph{Booleans, numerals, streams, codes.}  $\T = K$ and $\F = K\,I$; they are
normal and distinct, so by Church--Rosser no term reduces to both.  A boolean
selects: $\T\,x\,y \red x$ and $\F\,x\,y \red y$.  Scott numerals are
$\num{0} = K$ and $\num{m+1} = \fn{Succ}\,\num{m}$ with
$\fn{Succ} = S\,(K\,K)\,(S\,(K\,(S\,I))\,K)$, so that
$\fn{Succ}\,x\,z\,s \hred s\,(K\,x\,s) \red s\,x$.  Streams are Scott pairs:
$\fn{Cons}\,h\,t\,k \red k\,h\,t$, with $\fn{Hd}\,s = s\,\T$ and
$\fn{Tl}\,s = s\,\F$.  The code $\code{t}$ of a term is its Scott encoding over the
object type $S \mid K \mid I \mid \fn{App}$ with four continuations
(\lean{quoteT}), so $\code{f\,a} = \fn{App}\,\code{f}\,\code{a}$ literally.  For a
fixed code $c$, a term $\fn{EqT}_c$ decides equality with $c$:
$\fn{EqT}_c\,\code{d}\red(c = d)$.  We build it by recursion on $c$ in the
metalanguage (\lean{eqT\_spec}).

\paragraph{Recursion.}  Recursive terms are fixed points of a step function $F$,
written $\fn{FixSelf}\,F\,z$ after \texttt{lean-ski}.  Each unrolling wraps the
self-reference in one more $I$, so the self-references are not syntactically equal;
the invariant is that each still head-reduces to $\fn{FixB}\,F$
(\lean{fixSelf\_unroll}).

\paragraph{Infinite numerals.}  Let $\omega$ range over the self-references
$\fn{FixSelf}\,\fn{Succ}\,z$ with $z \hred \fn{FixB}\,\fn{Succ}$:
the \emph{infinite numerals}.  With two arguments,
$\omega\,a\,s \hred s\,(K\,\omega'\,s)$ for another infinite numeral $\omega'$, in
one explicit step followed by a head reduction (\lean{omega\_steps}), exactly as
$\num{m+1}\,a\,s\hred s\,(K\,\num{m}\,s)$; with fewer arguments $\omega$ reaches a
weak head normal form other than $K$ (\lean{omega\_whnf0}, \lean{omega\_whnf1}).

\paragraph{Prefixes and layers.}  A term $s$ \emph{has the prefix}
$w = b_0\cdots b_{k-1}$ when $s \red \fn{Cons}\,b_0\,s_1$,
$s_1 \red \fn{Cons}\,b_1\,s_2$, \dots, $s_{k-1}\red \fn{Cons}\,b_{k-1}\,s_k$
for some $s_i$, the last tail arbitrary (\lean{HasPrefix}).  A term $n$ \emph{has at
least $m$ layers} when $n \red \fn{Succ}\,n_1$, $n_1\red\fn{Succ}\,n_2$,
\dots, $m$ times, the last tail arbitrary (\lean{AtLeast}).  Both are closed under
backward reduction, and $\num{c}$ has at least $m$ layers when $c \ge m$.

%% ---------------------------------------------------------------------------
\section{Staging and behavioral lifts}\label{sec:staging}
Compile-time quotation makes $\Box$ a comonad without a unit (\S\ref{sec:modality}).
A behavioral lift is a weak substitute for the missing unit (\S\ref{sec:lift}).
\subsection{The staging modality}\label{sec:modality}
S4 is the modal logic with the axioms T, $\Box A\to A$ (what is necessary is true),
and 4, $\Box A\to\Box\Box A$ (if something is necessary, it is necessarily necessary).
Davies and Pfenning read $\Box A$ as codes of terms of type $A$
\citep{DaviesPfenning2001}: T is then evaluation of code, and 4 re-encodes a code as
the code of that code.  The converse of T, $A\to\Box A$, which S4 does not prove,
would quote a live value.  Over $\{S,K,I\}$ the first two are definable and the
third is not (Table~\ref{tab:s4}).
\begin{table}[htb]\centering\footnotesize
\caption{The S4 reading of staging, and what $\{S,K,I\}$ defines.}\label{tab:s4}
\begin{tabular}{@{}lll@{}}
\toprule
axiom & map & over $\{S,K,I\}$ \\
\midrule
T: $\Box A\to A$ & unquote, $\varepsilon$ & $\varepsilon\,\code{t}\red t$, definable \\
4: $\Box A\to\Box\Box A$ & re-encode, $\delta$ & $\delta\,\code{t}\red\code{\code{t}}$, definable (Prop.~\ref{prop:requote}) \\
converse of T: $A\to\Box A$ & quote, $\eta$ & no closed $Q$ has $Q\,M\red\code{M}$ \citep{Barendregt1991} \\
\bottomrule
\end{tabular}
\end{table}

An unquoter
$\varepsilon : \Box A \to A$ with $\varepsilon\,\code{t} \red t$ is definable
(e.g., SKIjack encodes one in 43
atoms \citep{DavisMeta}).  Re-encoding is definable as well since
it acts on data: \newpage
\begin{proposition}[\lean{Requote.requote}]\label{prop:requote}
There is a closed $R$ with $R\,\code{t} \red \code{\code{t}}$ for every $t$.
\end{proposition}
\begin{proof}
$R$ is the fixed point of
\[
\lambda r\,q.\ q\,\code{\code{S}}\,\code{\code{K}}\,\code{\code{I}}\,
(\lambda f\,a.\ \fn{App}'\,(r\,f)\,(r\,a)),
\]
where
$\fn{App}'\,u\,v = \fn{App}\,(\fn{App}\,\code{\fn{App}}\,u)\,v$ is the
code of an application given the codes of its parts.  Induction on $t$.
\end{proof}
So $\varepsilon$ and $\delta = R$ give $\Box$ the structure of a comonad, and the
unit $\eta : A \to \Box A$, a closed $Q$ with $Q\,M\red\code{M}$, is exactly what no
term provides \citep{Barendregt1991} (\lean{no\_quote}).  Kavvos
\citep[\S1.3--1.4]{Kavvos2017} describes the same situation for untyped $\lambda$.
Self-reference survives the missing unit:
\begin{theorem}[Intensional fixed points; \lean{Ifix.ifix}, \lean{Ifix.quine}]
For every $g$ there is $t$ with $t \red g\,\code{t}$.  With $g = I$, $t\red\code{t}$.
\end{theorem}
\begin{proof}
Let $h = \lambda y.\ g\,(\fn{App}\,y\,(R\,y))$ and $t = h\,\code{h}$.  Then
\[
t \red g\,(\fn{App}\,\code{h}\,(R\,\code{h})) \red
g\,(\fn{App}\,\code{h}\,\code{\code{h}}) = g\,\code{h\,\code{h}} = g\,\code{t}.
\]
The only quotation is the compile-time code of the closed term $h$.
\end{proof}
This is Kavvos' intensional fixed-point rule \citep{Kavvos2017}, realized by
re-encoding.

\subsection{Behavioral lifts}\label{sec:lift}
\begin{definition}[Lift]
Fix a set $X$ of test inputs and a hypothesis $H$ on terms.  A closed term $L$
\emph{lifts} under $H$ on $X$ when for every $f$ with $H(f)$ there is a term $c$ with
\[
L\,f \red \code{c} \qquad\text{and}\qquad
\forall x\in X,\ \forall r\in\{\T,\F\},\ \ f\,x\red r \implies c\,x \red r.
\]
\end{definition}
The definition is as weak as the informal notion allows: $L$ need work only on terms
satisfying $H$; $c$ may be any term and need be neither $f$ nor canonical; agreement
is required only on $X$.
The impossibility (\S\ref{sec:nolift}) refutes even the weaker form in which
$L\,f$ and $\code{c}$ need only have a common reduct.  The impossibility is not a corollary of the
no-quote theorem, which concerns syntax and every term, total or not; the finite lift
below exists although that theorem holds.

The instances are listed in Table~\ref{tab:results}, with the test inputs for
$\Bool\to\Bool$ as $\T,\F$ and for $\Nat\to\Bool$ as the numerals.  They are, at Cantor
space, the terms with a prefix of length $m$ under a modulus, and the definable
streams under totality.

A \emph{Q-coalgebra} in Kavvos' sense \citep[Def.~33]{Kavvos2017} is a map
$\alpha : A\to\Box A$ with $\varepsilon\circ\alpha\sim\mathrm{id}$ and
$\delta\circ\alpha\sim\Box\alpha\circ\alpha$.  The first law, with $\sim$ read as
agreement on the test inputs, is satisfied by the lifts of \S\ref{sec:modulus}.
The second law is discussed in \S\ref{sec:coalg}.

%% ---------------------------------------------------------------------------
\section{Lifts from a modulus}\label{sec:modulus}
Each lift in this section has the same shape: a search decides properties of the
function's finite table, a tabulation turns the table into a canonical tree, and the
lift returns the code of an evaluator for that tree.  The finite case needs no search,
since its domain is listed.

\subsection{The finite case}\label{sec:bool}
\begin{theorem}[\lean{LiftBool.lift\_bool}, \lean{lift\_canonical}]\label{thm:bool}
The term
\[
L_{\Bool} = \lambda f.\ f\,\T\,(f\,\F\,\code{K\,K}\,\code{I})\,
(f\,\F\,\code{\fn{Not}}\,\code{K\,(K\,I)})
\]
lifts under totality on $\{\T,\F\}$.
Functions with the same values get the same code.
\end{theorem}
\begin{proof}
If $f\,\T\red t$ and $f\,\F\red e$, two selections give the code of the canonical
function with values $t, e$: $K\,K$ (both true), $I$, $\fn{Not}$, or $K\,(K\,I)$
(both false).  Each agrees with $f$ on $\T$ and $\F$, because a term reaches at
most one boolean.
\end{proof}
This is routine; \citet{AltenkirchUustalu2004} define the same tabulation at every
finite type, in the metalanguage.

\subsection{At \texorpdfstring{$\Nat\to\Bool$}{Nat -> Bool}}\label{sec:nat}
\begin{definition}[$m$-local; \lean{NatLift.NLocal}]
$f$ is \emph{$m$-local} with values $\varphi(0),\dots,\varphi(m)$ when
$f\,\num{c}\red\varphi(c)$ for $c<m$, and $f\,n\red\varphi(m)$ for every term $n$ with
at least $m$ layers.
\end{definition}
An $m$-local $f$ forces at most $m$ successor layers of its argument, and $f$ has a
\emph{modulus} when it is $m$-local for some $m$.  An $m$-local $f$ is total on the
numerals (\lean{nlocal\_total}) and reaches $\varphi(m)$ on an infinite numeral
(\lean{omega\_atLeast}), so it is total on the points of $\mathbb{N}_\infty$, with
value $\varphi(m)$ at $\infty$. 

\paragraph{Search.}  Let
$\fn{Srch}\,p = p\,\num{0}\,\num{0}\,(\fn{Succ}\,(\fn{Srch}\,(p\circ\fn{Succ})))$:
if $p\,\num 0$ then $\num 0$, else the successor of the search for
$p\circ\fn{Succ}$.  It builds a lazy numeral, infinite when $p$ is never true.
\begin{lemma}[Search over $\mathbb{N}_\infty$; \lean{search\_nat}]\label{lem:searchN}
If $p$ is $m$-local with values $\varphi$, then $p\,(\fn{Srch}\,p)$ reduces to
$\T$ iff $\varphi(k)$ holds for some $k\le m$, and to $\F$ otherwise.
\end{lemma}
\begin{proof}
Induction on $m$.  For $m=0$, $p$ is constant.  For $m+1$, $p\circ\fn{Succ}$ is
$m$-local with values $\varphi(\cdot+1)$ (\lean{nlocal\_shift}).  If $\varphi(0)$,
$\fn{Srch}\,p\red\num 0$ and $p\,\num 0\red\T$.  Otherwise
$\fn{Srch}\,p\red\fn{Succ}\,(\fn{Srch}\,(p\circ\fn{Succ}))$, and
$p\,(\fn{Succ}\,x)$ and $(p\circ\fn{Succ})\,x$ have a common reduct, so the
induction hypothesis and Church--Rosser give the result.
\end{proof}
This search corresponds to Escard\'o's for $\mathbb{N}_\infty$ \citep{Escardo2008}
but run on Scott numerals. $m$-locality lets $p$ terminate on the lazy numeral built
by the search.

\paragraph{Tabulation.}  The search yields two deciders.  The first,
$\fn{Every}\,p = \fn{Not}\,(p'\,(\fn{Srch}\,p'))$ with
$p' = \fn{Not}\circ p$, decides whether $\varphi$ holds at every $k\le m$
(\lean{everyN\_spec}).  The second,
$\fn{Const}\,f = \fn{Every}\,(\lambda n.\ \fn{EqB}\,(f\,n)\,(f\,\omega))$,
decides whether $f$ is constant (\lean{constN\_spec}).  Then
\[
\fn{Tree}\,f = \fn{Const}\,f\ (\fn{Leaf}\,(f\,\omega))\ 
(\fn{Node}\,(f\,\num 0)\,(\fn{Tree}\,(f\circ\fn{Succ})))
\]
reduces to a datum $D_m(\varphi)$ defined by recursion on $m$: a leaf of $\varphi(m)$ if
$\varphi$ is constant on $0..m$, else a node of $\varphi(0)$ and $D_{m-1}(\varphi(\cdot+1))$
(\lean{treeN\_spec}).  It depends only on $\varphi$ (\lean{DN\_congr}).
$\fn{Eval}\,d\,n = d\,(\lambda b.\,b)\,(\lambda b\,t.\ n\,b\,(\lambda n'.\ \fn{Eval}\,t\,n'))$
runs the datum: $\fn{Eval}\,D_m(\varphi)\,\num c\red\varphi(c)$ for $c<m$ and
$\fn{Eval}\,D_m(\varphi)\,n\red\varphi(m)$ for $n$ with at least $m$ layers
(\lean{evalN\_spec}).  The structural lift $\fn{LiftTree}$ at the tree's data type
sends the datum to its own code
(\lean{liftTreeN\_spec}).

\begin{theorem}[\lean{NatLift.nat\_lift}, \lean{nat\_canonical}]\label{thm:nat}
The term
\[
L_{\Nat} = \lambda f.\ \fn{App}\,\code{\fn{Eval}}\,(\fn{LiftTree}\,(\fn{Tree}\,f))
\]
lifts every $m$-local $f$, for every $m$, on the numerals:
$L_{\Nat}\,f \red \code{\fn{Eval}\,D_m(\varphi)}$, and that term agrees with $f$ on
every numeral.  Functions with the same values get the same code.
\end{theorem}
\begin{proof}
Compose the lemmas; agreement on $\num c$ with $c\ge m$ uses that $\num c$ has at least
$m$ layers, and that a term reaches at most one boolean.
\end{proof}

\subsection{At Cantor space}\label{sec:cantor}
\begin{definition}[$m$-local, streams; \lean{Search.Local}]
$q$ is \emph{$m$-local} with table $\varphi : \{\T,\F\}^m\to\{\T,\F\}$ when
$q\,s \red \varphi(w)$ for every term $s$ with a prefix $w$ of length $m$.
\end{definition}
As at $\Nat\to\Bool$, $q$ has a \emph{modulus} when it is $m$-local for some $m$.
\paragraph{Berger's search.}
The search \citep{Berger1990} is
\[
\fn{Find}\,p = (\lambda l.\ p\,l\,l\,(\fn{Cons}\,\T\,(\fn{Find}\,(p\circ\fn{Cons}\,\T))))\,
(\fn{Cons}\,\F\,(\fn{Find}\,(p\circ\fn{Cons}\,\F))).
\]
\begin{lemma}[\lean{Search.search}, \lean{search\_decides}]\label{lem:search}
If $p$ is $k$-local with table $\varphi$, then $\fn{Find}\,p$ has a prefix $u$ of
length $k$, $p\,(\fn{Find}\,p)\red\varphi(u)$, and $\varphi(u)$ holds whenever
$\varphi$ holds at some word of length $k$.
\end{lemma}
\begin{proof}
Induction on $k$.  For $k=0$, $p$ is constant, so $\fn{Find}\,p$ reduces to one of
the two cons cells.  For $k+1$, $p\circ\fn{Cons}\,b$ is $k$-local with table
$\varphi(b\,\cdot)$ (\lean{local\_pre}).  By induction $\fn{Find}\,(p\circ\fn{Cons}\,\F)$
has a prefix $u_0$, so $l = \fn{Cons}\,\F\,(\cdots)$ has the prefix $\F u_0$, and
$p\,l\red\varphi(\F u_0)$.  If that is $\T$, $\fn{Find}\,p\red l$.  Otherwise
$\varphi(\F\,\cdot)$ is constantly $\F$ (by the induction hypothesis for the left
branch), and $\fn{Find}\,p$ reduces to the right cons cell, whose prefix $\T u_1$
satisfies $\varphi$ if anything starting with $\T$ does.
\end{proof}
We claim no termination of $\fn{Find}$ but do not require it; the lemma asserts
only that reductions exist.  With sharing, $\fn{Find}\,p$ calls $p$ at most
$2^{k+1}-1$ times, and a predicate that reads its bits last-first comes close to
that bound.

\paragraph{Tabulation and lift.}  As in \S\ref{sec:nat}: $\fn{Every}$ and
$\fn{Const}$ (against the zero stream $0^\omega = \fn{Tab}\,(K\,\F)$) decide universality
and constancy (\lean{every\_spec}, \lean{const\_spec});
\[
\fn{Tree}\,q = \fn{Const}\,q\,(\fn{Leaf}\,(q\,0^\omega))\,
(\fn{Node}\,(\fn{Tree}\,(q\circ\fn{Cons}\,\F))\,(\fn{Tree}\,(q\circ\fn{Cons}\,\T)))
\]
reduces to a datum $D_m(\varphi)$ depending only on $\varphi$ (\lean{tree\_spec},
\lean{D\_congr}); an evaluator reads the stream's bits down the tree
(\lean{eval\_spec}); and $\fn{LiftTree}$ sends the tree to its code (\lean{liftTree\_spec}).

\begin{theorem}[\lean{CantorLift.cantor\_lift}, \lean{cantor\_canonical}, \lean{fun\_lift}]\label{thm:cantor}
The term
\[
L_C = \lambda q.\ \fn{App}\,\code{\fn{Eval}}\,(\fn{LiftTree}\,(\fn{Tree}\,q))
\]
lifts every $m$-local stream functional, for every $m$, on the terms with a prefix
of length $m$.  Functionals with the same table get the same code.  The same holds for functionals of functions
on numerals that read their argument at $0,\dots,m-1$ only, through
$\fn{Nth}$ and $\fn{Tab}$ (\lean{flocal\_stream}, \lean{tab\_spec}).
\end{theorem}
While the search and the tabulation are due to Berger and Escard\'o
\citep{Berger1990,Escardo2008,EscardoBlog2007}, the short step from them to a finite
canonical code is, as far as we know, novel.

\subsection{The second coalgebra law}\label{sec:coalg}
Our lifts satisfy Kavvos' second law up to
conversion of the quoted terms, but not with his $\sim$ at $\Box\Box A$.  For a lift
$\alpha$ with $\alpha\,f\red\code{c}$, $\delta(\alpha\,f)\red\code{\code{c}}$ and
$(\Box\alpha)(\alpha\,f) = \fn{App}\,\code{\alpha}\,(\alpha\,f)\red\code{\alpha\,c}$.  Each of our lifts is
idempotent on its outputs, $\alpha\,c\red\code{c}$: the output term meets the lift's own
hypothesis with the same table, and canonicality returns the same code
(\lean{Coalgebra.bool\_idem}, \lean{nat\_idem}, \lean{cantor\_idem}).  By
idempotence the two sides of the law are codes of convertible terms
(\lean{Coalgebra.coalg2}).  Kavvos' $\sim$ at $\Box\Box A$ is stricter.  An object is
\emph{point-univalent} when its extensionally equal points are intensionally equal
\citep[Def.~26]{Kavvos2017}, and every object of codes is
\citep[Lemma~18]{Kavvos2017}, so at $\Box\Box A$ his $\sim$ compares codes as codes.
The law fails there: $\code{\code{c}}$ and $\code{\alpha\,c}$ are codes of distinct
terms, and the equality test $\fn{EqT}$ of \S\ref{sec:prelim} separates them.  His results on Q-coalgebras, which assume the
law at that $\sim$ \citep[Lemmas~19--20, Thms.~30--31]{Kavvos2017}, therefore do not
apply to our lifts.

%% ---------------------------------------------------------------------------
\section{Totality and compactness}\label{sec:fan}
Section~\ref{sec:modulus} lifts every function with a modulus.  This section asks when
totality gives a modulus.  On the numerals it does not, and no lift exists
(\S\ref{sec:nolift}, \S\ref{sec:sharp}).  On the points of $\mathbb{N}_\infty$ it does
outright (\S\ref{sec:natinf}).  At Cantor space it does not on the definable streams,
and it does on all points in an oracle semantics (\S\ref{sec:fancantor}).  One simulation argument,
Lemma~\ref{lem:sim} and its generalizations, drives every cell.

\subsection{No lift under totality on the numerals}\label{sec:nolift}
\begin{theorem}[\lean{NoLift.no\_lift\_conv}]\label{thm:nolift}
No closed $L$ sends every $f$ total on the numerals to a term with a common reduct
with a code $\code{c}$ such that $c$ agrees with $f$ on every numeral.
\end{theorem}

\paragraph{Ingredients.}  Let $G\,c\,n = n\,\F\,(\lambda n'.\ c\,\T\,(\lambda c'.\ G\,c'\,n'))$, which decides
$n > c$ by walking down both numerals: $G\,\num{j}\,\num{n}\red(j<n)$
(\lean{G\_num}) and $G\,\omega\,\num{n}\red\F$ for all $n$ (\lean{G\_omega}).  The equality
test $\fn{EqT}_c$ of \S\ref{sec:prelim} suffices because the proof fixes $c$ first.

\begin{definition}
$\mathrm{Rel}_b\,t\,t'$ holds when $t'$ is $t$ with some infinite numerals replaced by
numerals $\num{m}$, $m\ge b$: it is the congruence generated by identity and the pairs
$(\omega,\num{m})$, $m\ge b$.
\end{definition}

\begin{lemma}[Stuttering simulation; \lean{NoLift.sim}]\label{lem:sim}
If $t \hred K$ in $n$ steps, $n < b$, and $\mathrm{Rel}_b\,t\,t'$, then $t'\hred K$.
\end{lemma}
The simulation is \emph{stuttering}: it matches runs, not single steps, as in the
simulation diagrams of compiler verification \citep{Leroy2009}.  At a
combinator head both sides contract the same redex.  At an infinite-numeral head,
$t$ unfolds $\omega$ and $t'$ unfolds $\num{m}$, by head reductions whose lengths may
differ, and the bound $b$ drops by one.  The bound counts how many more unfoldings of
$\omega$ the numerals $\num{m}$, $m\ge b$, can match.
\begin{proof}
Induction on $n$.  A related pair has a common shape: the same combinator at the head
of both spines with pointwise related arguments, or $\omega$ against $\num{m}$ with
related arguments (\lean{Rel.spine}).
\emph{Combinator heads.}  If the head is $S$, $K$ or $I$ with enough arguments, both
sides contract the same redex and the results are related at the same $b$ (a
congruence tolerates duplication by $S$ and erasure by $K$).  With too few arguments
$t$ is a weak head normal form, so $n=0$ and $t=K$, whence $t'=K$.
\emph{Infinite-numeral heads.}  With two or more arguments, $t$ takes at least one step
to $s\,(K\,\omega'\,s)\,\vec r$ and, since $m\ge b>n\ge 1$, $\num{m} = \num{k+1}$
reduces to $s'\,(K\,\num{k}\,s')\,\vec r\,'$.  The results are related at $b-1$, and
the remaining run from $t$'s reduct is shorter than $b-1$ steps.  With fewer
arguments $t$'s run would stop at a weak head normal form other than $K$, which
determinism of head reduction rules out.
\end{proof}

\begin{proof}[Proof of Theorem~\ref{thm:nolift}]
Suppose $L$ lifts.  The function $f = G\,\omega$ is total (constantly $\F$); let
$\code{c}$ be its code.  Then $\fn{EqT}_c\,(L\,f)\red\T$, so
$\fn{EqT}_c\,(L\,f)$ head-reduces to $K$ in some $n$ steps.  The term
$\fn{EqT}_c\,(L\,(G\,\num{n+1}))$ is related to it at $n+1$, so it too reduces to
$\T$ (Lemma~\ref{lem:sim}): $L$ gives the total function $G\,\num{n+1}$ the same
code $c$.  But $G\,\omega\,\num{n+2}\red\F$ while
$G\,\num{n+1}\,\num{n+2}\red\T$, so $c\,\num{n+2}$ reaches both booleans, contradicting
Church--Rosser.
\end{proof}
The argument is the classical one that the almost-everywhere-zero functions are not
finitely identifiable \citep{Balbach2020}: an answer committed to after finitely many
observations of the zero function is shared by a function that differs later.  The
simulation supplies ``finitely many observations'' for a live argument.

\subsection{Sharpness}\label{sec:sharp}
The proof of Theorem~\ref{thm:nolift} uses only $G\,\omega$ and the $G\,\num{k}$
(\lean{no\_lift\_core}), and each $G\,\num{k}$ has a modulus: it is $(k{+}1)$-local
(\lean{G\_num\_nlocal}).  So Theorems~\ref{thm:nolift} and~\ref{thm:nat} combine:
\begin{theorem}[\lean{G\_omega\_not\_local}, \lean{no\_lift\_local\_plus\_one}]\label{thm:sharp}
$G\,\omega$ is $m$-local for no $m$, and no closed term lifts every $m$-local function
(for every $m$) together with $G\,\omega$.
\end{theorem}
\begin{proof}
If $G\,\omega$ were $m$-local, $L_\Nat$ would lift it and every $G\,\num{k}$,
contradicting \lean{no\_lift\_core}; the second claim is the same argument for any $L$.
\end{proof}
So among the functions total on the numerals, the line between lift and no lift is
totality at $\infty$: adding $G\,\omega$, one function without a value there, defeats
every lift.

\subsection{At \texorpdfstring{$\mathbb{N}_\infty$}{N-infinity}}\label{sec:natinf}
Totality on the points of $\mathbb{N}_\infty$ gives a modulus outright.  The points of $\mathbb{N}_\infty$ are definable: the
numerals and an infinite numeral $\omega$.  Lemma~\ref{lem:sim} generalizes in two ways
(\lean{NatFan.simW}): an infinite numeral may be related to \emph{any} term with enough
successor layers, the right side reducing weakly; and the run may end at any weak head
normal form $v$, returning a related $v'$.  The second form serves the false case: a
run of $f\,\omega$ to $\F$ ends at some $K\,x$ with $x\red I$, and the simulation then
follows $x$'s run in turn.
\begin{theorem}[\lean{NatFan.nat\_fan}, \lean{local\_iff\_total\_inf}, \lean{nat\_lift\_inf}]\label{thm:natfan}
A function total on the numerals and at $\omega$ is $m$-local for some $m$.  Hence, for $f$
total on the numerals, $f$ has a modulus iff $f\,\omega$ reaches a boolean, and $L_\Nat$
lifts every $f$ total on the points of $\mathbb{N}_\infty$.
\end{theorem}
\begin{proof}
If $f\,\omega$ head-reduces to a boolean in $n$ steps (plus $n'$ for the inner $I$ in the
false case), the simulation fixes $f$'s value on every term with at least $n+n'+2$ layers.
The values on the finitely many numerals below that depth come from totality by finite choice, which needs no axiom of choice.
\end{proof}
With Theorem~\ref{thm:sharp}, the $\Nat$ row is settled as a statement about classes:
$L_\Nat$ lifts every function total on the numerals and at $\infty$, and no closed term
lifts that class together with $G\,\omega$, which is total on the numerals and has no value
at $\infty$.

\subsection{At Cantor space}\label{sec:fancantor}
At Cantor space the answer depends on which points count: totality on the definable
streams gives no lift, and totality on all points gives one in an oracle semantics.

\paragraph{Under totality on the definable streams.}  The impossibility of
\S\ref{sec:nolift} holds for any family $\Psi\,\omega$, $\Psi\,\num{k}$ that some test input
separates for every $k$ (\lean{CantorTotal.no\_lift\_family}).  Take a binary tree $T$ of
finite words, decided on stream prefixes by a closed term, and let $\Psi\,c\,s$ ask whether
$s$ leaves $T$ within $c$ steps.  $\Psi\,\num{k}$ gives up after $k$ steps;
$\Psi\,\omega$ never gives up.  If $T$ has no definable infinite path, $\Psi\,\omega$ is
total and true on every definable stream.  If $T$ is infinite and prefix-closed, a stream
that follows a long word of $T$ and then continues with zeros separates $\Psi\,\omega$
from $\Psi\,\num{k}$.
\begin{theorem}[\lean{KleeneTree.no\_lift\_cantor\_total'}]\label{thm:ctotal}
No closed term lifts every functional total on the definable streams.
\end{theorem}
The proof needs a Kleene tree decided by a closed term, and we build one inside the calculus
(\lean{KleeneTree.kleene}).  A small universal machine on codes performs one head step
(\lean{Univ.stepC\_spec}, mirroring the computable contraction) and a bounded run
(\lean{evalC\_spec}); a diagonal walk numbers every term (\lean{dec\_surj}, \lean{decC\_num}).
The tree holds a word $w$ iff, for every $e<|w|$, the term numbered $e$, applied to $\num{e}$
and read as a boolean, does not reach $w(e)$ within $|w|$ steps.  It is prefix-closed (a
verdict found with less fuel stands with more), infinite (choose each bit against its
verdict), and has no definable path: a definable stream $s$ gives the term
$\fn{Nth}\,s$, whose number $e$ has a run converging to $s$'s bit at $e$.

\paragraph{At Cantor space, with oracles.}  The step from pointwise totality to a modulus, the
\emph{fan step}, needs non-definable points.  Extend
the calculus with a constant $\fn{Ora}_\alpha$ for every $\alpha : \mathbb{N}\to\Bool$,
reducing as a stream cell: $\fn{Ora}_\alpha\,k \to k\,\alpha(0)\,\fn{Ora}_{\alpha(\cdot+1)}$.
We call the extended calculus the \emph{oracle semantics}: every point of Cantor
space, definable or not, is now a term.
A closed functional is \emph{oracle-total} when $F\,\fn{Ora}_\alpha$ head-reduces to
exactly $\T$ or $\F$ for every $\alpha$.  For $\F$ this asks more than the proof
needs; the relaxed form (\lean{OracleFan.cantor\_fan\_w}) uses the two-stage argument of
the $\mathbb{N}_\infty$ case.  A term may apply a stream to arbitrary continuations, not
only to the head and tail selectors, but the simulation absorbs this, because the oracle
answers any continuation as a cons cell would.  So a run of $F\,\fn{Ora}_\alpha$ is
mirrored on every term of the pure calculus that shares enough of $\alpha$'s bits,
whatever its tail (\lean{OracleFan.simO}).  Hence $F$ has a \emph{pointwise modulus}
at every $\alpha$: a length $n_\alpha$ such that $F$ returns its value at $\alpha$ on
every term that shares the first $n_\alpha$ bits of $\alpha$ (\lean{oracle\_local}).
\begin{theorem}[\lean{OracleFan.cantor\_fan}, \lean{cantor\_lift\_oracle}]\label{thm:ofan}
An oracle-total functional is $m$-local for some $m$, and $L_C$ lifts it.
\end{theorem}
\begin{proof}
The fan theorem (\lean{OracleFan.fan}, proved by K\"onig's lemma) turns the pointwise
moduli into one bound.  The fan theorem is classical: these results use
\texttt{Classical.choice}, and the simulation does not.
\end{proof}
\paragraph{Where the step is classical.}  Call $F$ \emph{barred} at a word $w$ when it is
constant on every term with prefix $w$, or when it is barred at both one-bit extensions
of $w$.  This is an inductive bar, the statement that $F$'s query tree is well-founded.
Constructively, being barred at the empty word is equivalent to having a modulus
(\lean{OracleBar.obar\_iff\_local}).  The fan step then factors into three parts
(\lean{cantor\_fan\_bar}):
\begin{itemize}
\item from runs to pointwise moduli, by the simulation (constructive);
\item from pointwise moduli to a bar, by the fan theorem for this bar
    (\lean{pointwise\_obar}; classical, the only such step);
\item from the bar to a table (constructive).
\end{itemize}

Passing to the term model does not remove the middle step; it makes the step fail.
Theorem~\ref{thm:ctotal} records the failure: over the definable streams, pointwise
totality gives no bar, because the Kleene tree above is infinite and has no definable
path.

\subsection{Scope}\label{sec:boundary}
Sections~\ref{sec:modulus} and~\ref{sec:fan} settle Table~\ref{tab:results} under one
standard hypothesis, the classical fan theorem in the oracle semantics.  An open question
in the $\Nat$ row remains: whether adding \emph{any} function without a value at
$\infty$ to the functions with a modulus defeats every lift.

%% ---------------------------------------------------------------------------
\section{Consequences for staged languages}\label{sec:skijack}
A staged language's data rule should refuse provable functions and admit lifted codes,
which are data.  A \emph{provable function} is a term that the checker can show, from its syntax and
the declared types, to be a function: a lambda, a macro, a core, a namespace literal,
or a combinator, constructor or equation applied to fewer arguments than it takes.
The guarantee is one-directional: it does not assert that every accepted occupant of a
data position is a datum \citep{skijack}.
SKIjack enforces the rule in two stages \citep{skijack}, the first of which is fixed
here.  Stage A (which formerly decided ``function'' by the presence of binders) decides
by arity, and refuses an under-applied combinator or equation, or a saturated equation
whose body is one.  Stage B, a Hindley--Milner discipline over the declared sum types,
refuses a function that reaches a data position through a bound variable.  With both
stages, no accepted program has a provable function in a data position.

Because lifted codes are data, $\mathtt{EQ}$ on two of them legitimately decides
extensional equality of two live functions at the liftable types.  The staging line is
not a confidentiality boundary: a term that holds a function can already probe it, and
a canonical code reveals only its extension.
The same line also appears one level up, in an interpreter hosted on a combinator calculus.
SKIlark \citep{skilark} implements Joy as SKIjack programs: its words are closed
supercombinators, and its quotations are Scott lists of words, so a Joy program is
data by construction and a word taken from a quotation is a symbol, not a live
function.  The no-quote theorem binds the host, not the object language: an
interpreter quotes its own programs because it represents them.  The staging line
reappears only where a live host value must become an object-level datum, which
is a lift in the sense of \S\ref{sec:lift}.

The arguments of \S\S\ref{sec:modulus}--\ref{sec:fan} use confluent weak reduction,
deterministic head reduction, and Scott encodings of booleans, numerals, streams and
codes.  Their case analyses are over $S$, $K$ and $I$, but analogous ones should go
through for any combinatory basis with these properties, each with its own
mechanization.  Calculi with an intensional operator differ: SF-calculus and tree
calculus \citep{JayGivenWilson2011,Jay2021} factorize normal forms, so the premise of
the no-quote argument changes there, and where the staging line may fall in them we
leave open.

%% ---------------------------------------------------------------------------
\section{Mechanization}\label{sec:mech}
All results are in \texttt{lean-ski} \citep{leanski}, on Lean~4.31.  The axiom gate \texttt{\#audit\_clean} admits only
\texttt{propext} and \texttt{Quot.sound}.  Every module's theorems are in the gate except
the fan theorem and the results that use it (\lean{OracleFan.fan}, \lean{cantor\_fan},
\lean{OracleBar.pointwise\_obar}), which add \texttt{Classical.choice}.  The table below
situates our results by module:
\begin{center}\small
\begin{tabular}{@{}ll@{}}
\toprule
module & results \\
\midrule
\lean{Requote}, \lean{Ifix} & re-encoding; intensional fixed points; a quine \\
\lean{LiftBool} & Theorem~\ref{thm:bool} \\
\lean{NatLift} & Lemma~\ref{lem:searchN}, Theorems~\ref{thm:nat} and~\ref{thm:sharp} \\
\lean{Search} & Lemma~\ref{lem:search}; $\fn{Nth}$, $\fn{Tab}$ \\
\lean{CantorLift} & Theorem~\ref{thm:cantor} \\
\lean{Coalgebra} & the second coalgebra law \\
\lean{NoLift} & Theorem~\ref{thm:nolift}: simulation, $G$, $\fn{EqT}$, \lean{no\_lift\_core} \\
\lean{NatFan} & Theorem~\ref{thm:natfan} \\
\lean{CantorTotal}, \lean{Univ}, \lean{KleeneTree} & Theorem~\ref{thm:ctotal} \\
\lean{OracleFan} & Theorem~\ref{thm:ofan} (\texttt{Classical.choice}) \\
\lean{OracleBar} & the classical step isolated \\
\bottomrule
\end{tabular}
\end{center}
We build recursive terms with \texttt{lean-ski}'s $\fn{FixSelf}$ and bracket
abstraction over open terms (\lean{absN}); all reasoning is by existential weak reduction, with head
reduction only inside the simulation.

%% ---------------------------------------------------------------------------
\section{Related work}\label{sec:related}
\paragraph{Modality and intensionality.}  S4 for staging is due to Davies--Pfenning
\citep{DaviesPfenning2001}.  Kavvos \citep{Kavvos2017,Kavvos2020} gave the intensional
comonad, Barendregt's no-quote argument, and intensional recursion with the
G\"odel--L\"ob axiom \citep{Kavvos2021}.  His Q-coalgebras \citep[Def.~33]{Kavvos2017}
are the abstract form of a behavioral lift; he did not characterize which types carry
one.  Another factorization impossibility for $SK$ is due to \citet{JayGivenWilson2011};
Jay's tree calculus quotes programs but not computations \citep{Jay2021}.
\paragraph{Reification.}  Normalization by evaluation
\citep{BergerSchwichtenberg1991} and type-directed partial evaluation
\citep{Danvy1996} reify at higher types by type-directed $\eta$-expansion over a
model with syntax at base type; sums need control or a sheaf model
\citep{Danvy1996,ADHS2001}, and inductive types are out of reach.
\citet{AltenkirchUustalu2004} quote at every finite type of $\lambda^{\to 2}$ by
enumerating finite domains as decision trees, in the metalanguage.  Our lifts are terms
of the calculus applied to live arguments.
\paragraph{Exhaustibility.}  Berger introduced the search over Cantor space
\citep{Berger1990}.  The fan functional computed from a functional alone, decidable
equality at compact domains, and the searchability of $\mathbb{N}_\infty$ are due to
Escard\'o \citep{Escardo2007,Escardo2008,EscardoBlog2007}; our reading of the staging line borrows
the shape of his compact and discrete types \citep[Thm.~4.13]{Escardo2008}.  Decision trees for
continuous functionals are also built from a term's syntax, by translation (dialogue
trees; see \citet{BaillonMahboubiPedrot2026}); ours come from the live value, by search.
\paragraph{Continuity.}  It is standard that a computation queries a type-1 argument finitely often,
and so cannot extract an index from an oracle \citep[\S2.3--2.4]{Longley2005};
see \citet{LongleyNormann2015}.  Our specific argument is finite identification
\citep{Balbach2020}.  We know of no continuity theorem for weak combinatory logic;
combinatory B\"ohm trees \citep{DybjerKuperberg2012} would be its natural home, and in
$\lambda$ the continuity and genericity lemmas \citep[\S14.3]{Barendregt1984} would give
our step.  The stuttering simulation is a self-contained substitute.


\bibliographystyle{plainnat}
\bibliography{mss}
\end{document}
